\documentclass[aspectratio=1610,12pt,notheorems]{beamer}

\usepackage[utf8x]{inputenc} \usepackage[russian]{babel}
\usepackage{amsmath,amssymb,amsthm,mathtools}
\usepackage{graphicx,caption,subcaption}
\usepackage{hyperref,natbib}
\usepackage{tikz,xcolor,makecell}
\usepackage{algorithm,algpseudocode}

\usetikzlibrary{arrows,backgrounds,patterns,%
	matrix,shapes,fit,calc,shadows,plotmarks,snakes}

\theoremstyle{plain}
\newtheorem{theorem}{Теорема}
\newtheorem{lemma}[theorem]{Лемма}

\theoremstyle{definition}
\newtheorem{definition}{Определение}
\newtheorem{problem}{Задача}

\usetheme[height=0.97cm]{Rochester}
\usecolortheme{dolphin}

\definecolor{hard}{RGB}{40,40,128} % {140,50,50}
\definecolor{mnsgold}{RGB}{255,255,130} % {240,230,200}

\setbeamercolor{headline}{bg=hard,fg=mnsgold}
\setbeamercolor*{frametitle}{parent=headline}

\setbeamercolor{structure}{fg=hard}
\setbeamercolor{subsection in head/foot}{bg=white,fg=hard}
\setbeamercolor{section in head/foot}{bg=hard,fg=mnsgold}
\setbeamercolor{block title}{bg=hard,fg=mnsgold}

\setbeamertemplate{navigation symbols}{}

\def\mitem{\medskip\item}
\def\ps{\\ [0.65cm]} \linespread{1.16} \parskip=2.6mm
\def\usl#1#2{\begin{block}{#1} #2 \end{block} \medskip\pause}

\def\ll{\left(} \def\rr{\right)}
\def\lag{\left\langle} \def\rag{\right\rangle}

\definecolor{failpos}{RGB}{230,30,20}
\definecolor{initpos}{RGB}{30,20,220}
\definecolor{turna}{RGB}{100,240,110}
\definecolor{turnb}{RGB}{250,140,110}


\title[Methodic keys to MNS]
    {\bfseries Методические приёмы \\
	составления задач олимпиады}

\author[\ ]
	{Б. А. Золотов,\quad «Математика НОН-СТОП»\\ \vspace{0.3cm}
		{\small Фонд «Время Науки»}}

\institute[\ ]{\ }

\date{8 ноября 2020}

%%%%%%%%%%%%%%
%%%%%%%%%%%%%%

\begin{document}

\frame{\titlepage}

\begin{frame}
	\ \\
	\noindent В 2019 году олимпиада «Математика НОН-СТОП» проводится
	с использованием гранта Президента Российской Федерации \\
	на развитие гражданского общества, предоставленного \\
	{\bfseries Фондом президентских грантов.}
	
	\vspace{0.45cm}
	\begin{center}
		\includegraphics[scale=0.6]{fpg}
	\end{center}
\end{frame}

\begin{frame}\frametitle{История олимпиады}
\begin{itemize}
	\item 2010 — первая олимпиада;
        %\mitem 2011 — 418 участников
%	\mitem 2012 — три площадки;
%	\mitem 2015 — первый профильный вариант;
	\item 2016 — 400 участников пишут базовый вариант, 92 --- профильный;\\
        $\phantom{2016}$ — поддержка Фонда «Время Науки»;
	%\mitem 2017 — задачи для 4 класса;\\
	%$\phantom{2017}$ — автоматическое распределение участников;
	\item 2018 — 847 участников пишут базовый вариант, 128 --- профильный;\\
        $\phantom{2018}$ — включение в Перечень региональных олимпиад и конкурсов\\
	$\phantom{2018 — }$\quad интеллектуальной направленности;\\
	$\phantom{2018}$ — поддержка Фонда Президентских грантов,\\
	$\phantom{2018 — }$\quad Комитета по образованию СПб;\\
	%$\phantom{2018}$ — этот семинар;\\
	\item 2019 — выход сборника задач;\\
        $\phantom{2019}$ — площадки в Бердске (Новосибирская обл.) и Гомеле (Беларусь);\\
        $\phantom{2019}$ — число участников приближается к 2000.
\end{itemize}\end{frame}

\begin{frame}\frametitle{Статистика олимпиады}
\begin{itemize}
	\item 12 площадок (на 2021 год — 20 соглашений);\\
	\item количество участников --- около 2000;\\
	\item пять городов: Санкт-Петербург, Бердск (Новосибирская обл.), \\
		Нов. Уренгой (ЯНАО), Гатчина (ЛО), \\
		Гомель (Беларусь);\\
	\item две страны;\\
	\item проблемы с часовыми поясами.
\end{itemize}\end{frame}

\begin{frame}\frametitle{Не верьте им!}
	Основная цель олимпиады — избегать появления «сообщества профессионалов», которые умеют её решать. \ps
	То, что вам сейчас расскажут, — не руководство по подготовке и не классификация задач будущей олимпиады; \ps
	А набор наблюдений, касающихся заданий прошлых лет.
\end{frame}


\section[Задачи на приведение примера]{1. Constructive Problems}

\begin{frame} \frametitle{Разрезания}
	\usl{2020-5-1C}{
		Можно ли нарисовать на клетчатом листе бумаги такую фигуру, которую можно разрезать по линиям сетки на две {\itshape одинаковые} фигуры двумя способами~— причём фигуры в первом и во втором способе были бы одни и те же, но линии разреза выглядели бы по-разному?
	}

\begin{center} \tikz[scale=0.75]
{
   \foreach \x / \y in {
	1/3, 2/3, 3/3, 3/4, 6/1, 6/2, 6/3, 7/3
   } 
	\draw[thick,fill opacity=0.6] (\x, \y) rectangle (\x cm + 1 cm, \y cm + 1 cm);
   {
   \foreach \x / \y in {
	1/1, 1/2, 2/2, 3/2, 7/2, 8/2, 8/3, 8/4
   } 
	\draw[thick,fill=gray,fill opacity=0.6] (\x, \y) rectangle (\x cm + 1 cm, \y cm + 1 cm);
	\node at (1.5,2.5) {$A$};
	\node at (3.5,3.5) {$B$};
	\node at (6.5,3.5) {$B$};
	\node at (8.5,2.5) {$A$};
}
} 
\end{center}
\end{frame}

\begin{frame} \frametitle{Сумма цифр}
\usl{2017-8-2C}{
	Придумайте (или расскажите, как построить) 95-значное число, в котором нет нулей и которое делится на свою сумму цифр.
}
Придумаем число, делящееся на $144=9\cdot 16$:
$$\underbrace{.\ .\ .\ .\ .\ .\ .}_{\begin{minipage}{2cm}
	\scriptsize разрядов — 91, \\
	$\sum$ цифр — 134
\end{minipage}}3232.$$
\end{frame}

\begin{frame}
\frametitle{Лабиринт}

\usl{2019-7-8C}{
Путник в лабиринте видит ситуацию вокруг. Помимо этого, никакой другой 
информации и памяти у него нет. Существует ли какой-нибудь набор правил, чтобы он, 
имея только эту информацию, мог дойти до финальной клетки в любом лабиринте?} \vspace{4mm}

Заметим, что поведение путника однозначно определено в простых ситуациях:

\begin{center} \tikz[scale=0.8]{\begin{scope}[xshift=-3cm]
	\fill[pattern=north east lines] (0,0) -- (0,1) -- (1,1) -- (1,0) -- (1.7,0) -- (1.7,1.7) -- (-0.7,1.7) -- (-0.7,0) -- cycle;
	\draw (0,0) -- (0,1) -- (1,1) -- (1,0);
	\filldraw (0.5,0.5) circle[radius=1.2mm]; \end{scope}
\begin{scope}[xshift=2cm]
	\fill[pattern=north east lines] (-0.7,-0.2) rectangle (0,1.2) (1,-0.2) rectangle (1.7,1.2);
	\draw (0,-0.2) -- (0,1.2) (1,-0.2) -- (1,1.2);
	\filldraw (0.5,0.5) circle[radius=1.2mm];
\end{scope}} \end{center}
\end{frame}

\begin{frame}
\frametitle{Лабиринт}

Приведём решение без $T$-образных перекрёстков, чтобы о них не думать: \ps\pause

\begin{center} \tikz[scale=1.14]{
	\draw (-3,0) ++(0.15,0.15) rectangle ++(5.7,3.7);
	\foreach \x in {-3,...,2} {
	    \foreach \y in {0,...,3} {
		\draw[pattern=north east lines] (\x, \y) ++(0.15,0.15) --
		    ++(0,0.7) -- ++(0.7,0) -- ++(0,-0.7) -- ++(-0.7,0);
	    };
	};
	\draw (2,3.65) node{\itshape\footnotesize Ф};
	\filldraw (0,1.5) circle[radius=0.56mm];
} \end{center} \end{frame}


\section[Задачи, вдохновлённые реальностью]{2. Reality-Inspired Problems}

\begin{frame} \frametitle{Семнадцатый независимый}

\usl{2020-5-3A}{
	Песню каждого участника оценивает 15 судей. Судья ставит каждому участнику в паре от 0 до 22 баллов и отдаёт свой голос участнику, которому поставил больше баллов. В паре объявляется победителем тот участник, которому отдано больше голосов. Может ли быть так, что победитель в паре набрал меньше баллов, чем проигравший, несмотря на перевес в голосах?
}

\end{frame}

\begin{frame} \frametitle{Семнадцатый независимый}

\newcommand{\vs}{&\hspace{-3mm}:&\hspace{-3mm}}
\begin{center} \begin{tabular}{|l|rcl|rcl|c|c|}
	\hline
	Участник & \multicolumn{3}{c|}{Победы}
		& \multicolumn{3}{c|}{Проигрыши}
		& Баллы & Голоса \\ \hline \hline
	Победитель 1 & 1 \vs 0 & 0 \vs 22 & 11 & 11 \\ \hline
	Победитель 2 & 2 \vs 0 & 0 \vs 16 & 20 & 10 \\ \hline
	Победитель 3 & 3 \vs 0 & 0 \vs 12 & 27 & 9 \\ \hline
	Победитель 4 & 4 \vs 0 & 0 \vs 9 & 32 & 8 \\ \hline \hline
	Проигравший 4 & 9 \vs 0 & 0 \vs 4 & 63 & 7 \\ \hline
	Проигравший 3 & 12 \vs 0 & 0 \vs 3 & 72 & 6 \\ \hline
	Проигравший 2 & 16 \vs 0 & 0 \vs 2 & 80 & 5 \\ \hline
	Проигравший 1 & 22 \vs 0 & 0 \vs 1 & 88 & 4 \\ \hline
\end{tabular} \end{center}

\end{frame}

\begin{frame}
\frametitle{Аксиомы выборов}

\usl{2019-7-3A}{На предприятии работают 50 человек, и они выбирают себе начальника. Есть две кандидатуры, Ваня и Даня. Про каждого работника известно заранее, кому он отдаёт предпочтение: 20 человек за Даню, 30 человек за Ваню. \smallskip \\
Голосование проходит по двухтуровой системе: люди делятся на 5 групп по 10 человек, в каждой группе выбирается кандидат, наиболее популярный среди членов этой группы, и затем из 5 ответов выбирается имя, названное большее число раз. \smallskip \\
Разделите работников на группы так, чтобы в большинстве групп выбрали Даню и он победил на выборах, несмотря на изначально меньшее число голосующих за него.}
\end{frame}

\begin{frame}
\frametitle{Аксиомы выборов}

\begin{center} \tikz{
\begin{scope}[rotate=90,scale=1.2]
	\filldraw[fill=gray,draw=gray,opacity=0.32] (0,0) rectangle (1,5);
	\foreach \x in {0,2,5} {\draw[thick, color=gray]
		(0.5 * \x cm, 0) -- (0.5 * \x cm, 5);}
	\foreach \x in {0,10} {\draw[thick, color=gray]
		(0, 0.5 * \x cm) -- (2.5, 0.5 * \x cm);}
	\foreach \x in {1,3,4} {\draw[color=gray, opacity=0.38]
		(0.5 * \x cm, 0) -- (0.5 * \x cm, 5);}
	\foreach \x in {1,...,9} {\draw[color=gray, opacity=0.38]
		(0, 0.5 * \x cm) -- (2.5, 0.5 * \x cm);}

	\draw (0.5,-0.25) node[right]{\large За Даню};
	\draw (1.75,-0.25) node[right]{\large За Ваню};

	\pause
	\begin{scope} [xscale=-1, xshift=-2.5cm]
	\draw[very thick] (0.5,0) -- (0.5,5);
	\draw[very thick] (2.5,1.5) -- ++ (-1,0) --
		++ (0,0.5) -- ++ (-0.5,0) -- ++(0,-2);
	\draw[very thick] (2.5,3.5) -- ++ (-1,0) --
		++ (0,-0.5) -- ++ (-0.5,0) -- ++(0,2);
	\draw[very thick] (2.5,3.5) --
		++ (-1,0) -- ++ (0,-0.5) --
		++ (-0.5,0) -- ++ (0,-1) -- ++ (0.5,0) --
		++ (0,-0.5) -- ++(1,0);
	\end{scope}
\end{scope} } \end{center} \end{frame}

\begin{frame} \frametitle{2020-4-1A, Эскалаторы}

\begin{center} \begin{tabular}{cc}
	\makecell[l]{\newcommand{\staircs}{\filldraw[draw=black,fill=white,fill opacity=0.84]}
\newcommand{\stairwd}{0.305}

\newcommand{\flr}[1]{
     \staircs (0, #1 cm)
	-- ++(-0.5 * \stairwd cm, -1 * \stairwd cm) -- ++(5.8,0)
	-- ++(0.5 * \stairwd cm, \stairwd cm) -- ++(1,2)
	-- ++(0.5 * \stairwd cm, \stairwd cm) -- ++(-5.8,0)
	-- cycle;
}

\newcommand{\stairst}[1]{ \begin{scope}[yshift = #1 cm]
     \staircs (1,2)
	-- ++(0.5 * \stairwd cm, \stairwd cm) -- ++(5.8,4)
	-- ++(-1 * \stairwd cm, -2 * \stairwd cm) -- ++(-5.8,-4)
	-- cycle;
     \staircs (6.3,1)
	-- ++(0.5 * \stairwd cm, \stairwd cm) -- ++(-5.8,4)
	-- ++(-1 * \stairwd cm, -2 * \stairwd cm) -- ++(5.8,-4)
	-- cycle;
     \staircs (0,0)
	-- ++(0.5 * \stairwd cm, \stairwd cm) -- ++(5.8,4)
	-- ++(-1 * \stairwd cm, -2 * \stairwd cm) -- ++(-5.8,-4)
	-- cycle;
     \flr{4}
\end{scope}}

\tikz[scale=0.34]{
	\flr{0} \stairst{0} \stairst{4}
	\stairst{8} \stairst{12}
	\draw[<->] (-0.5 * \stairwd cm, -0.9) -- (5.8 cm - 0.5 * \stairwd cm,-0.9);
	\draw (2.9 cm - 0.5 * \stairwd cm, -0.9) node[below]{12\,м};
	\draw[<->] (5.8 cm + 0.9 cm - 1 * \stairwd cm, -1 * \stairwd cm)
	     -- (6.8 cm + 0.9 cm, 2 cm + \stairwd cm);
	\draw (7.2 cm - 0.5 * \stairwd cm, 1 cm) node[right]{4\,м};
}} &
	\makecell[l]{
		В торговом центре Атмосфера» пять этажей. \\
		Соседние этажи соединяют по три эскалатора. \\
		По утрам администратор запускает каждый \\
		из эскалаторов вверх или вниз. Cпособ \\
		запустить эскалаторы {\it хороший,} если \\
		с любого этажа можно добраться до любого, \\
		не проходя ни по одному из этажей \\
		более 5 метров. Приведите пример; сколько всего \\
		хороших способов ? ($2 \cdot 3^4 = 162$.) \\ \ \\ \ \\}
\end{tabular} \end{center}

\end{frame}


\section[Задачи на приведение множества примеров]{3. Many Examples}

\def\dwtc{\draw[very thick] }
\def\dtc{\draw[thick] }
\def\dtcg{\draw[thick,color=gray,opacity=0.68] }
\def\dcw{\draw[color=white] }

\def\hexnhex#1{\tikz[scale=0.73]{
    \dcw (-0.3,1) rectangle (0,2); \dcw (3,1) rectangle (3.3,2);
    \foreach \x in {1,...,5} {
	\draw[color=gray, opacity=0.38] (0, 0.5 * \x cm) -- (3, 0.5 * \x cm);
	\draw[color=gray, opacity=0.38] (0.5 * \x cm, 0) -- (0.5 * \x cm, 3);
    }
    \dtc (0,0) -- (0,3) -- (3,3) -- (3,0) -- cycle;
    \dtc #1;
}}

\def\kra{-- ++(0,-0.5) -- ++(-0.5,0) -- ++(0,-0.5) -- ++(0.5,0)}
\def\krb{-- ++(0,-0.5) -- ++(0.5,0) -- ++(0,-0.5) -- ++(-0.5,0)}
\def\tabrow#1{\begin{center}\begin{tabular}{ccc} #1 \end{tabular}\end{center}}

\begin{frame}
\frametitle{Разрезай и властвуй}
\usl{2019-7-6C}{Предложите как можно больше разных способов разрезать квадрат $6 \times 6$ на два одинаковых многоугольника по линиям сетки.} \vspace{3mm}

\tabrow{
	\hexnhex{(1.5,0) -- (1.5,3)} &
	\hexnhex{(1,0) -- (1,1.5) -- (2,1.5) -- (2,3)} &
	\hexnhex{(0.5,0) -- (0.5,1.5) -- (2.5,1.5) -- (2.5,3)}}
\end{frame}

\begin{frame} \frametitle{Разрезай и властвуй}
\tabrow{
	\hexnhex{(1,0) -- (1,1) -- (1.5,1) -- (1.5,2) -- (2,2) -- (2,3)} &
	\hexnhex{(0.5,0) -- (0.5,1) -- (1.5,1) -- (1.5,2) -- (2.5,2) -- (2.5,3)} &
	\hexnhex{(0.5,3) -- (0.5,1) -- (1.5,1) -- (1.5,2) -- (2.5,2) -- (2.5,0)}}
\tabrow{
	\hexnhex{(1,3) -- (1,0.5) -- (1.5,0.5) -- (1.5,2.5) -- (2,2.5) -- (2,0)} &
	\hexnhex{(1,3) \kra \kra -- (1,0.5) -- (1.5,0.5)
	    -- (1.5,2.5) -- (2,2.5) \krb \krb -- (2,0)} &	
	\hexnhex{(1,3) -- (1,1) -- (1.5,1) -- (1.5,2) -- (2,2) -- (2,0)}}
\tabrow{
	\hexnhex{(1,3) -- (1,1) -- (0.5,1) -- (0.5,0.5) -- (1.5,0.5)
	    -- (1.5,2.5) -- (2.5,2.5) -- (2.5,2) -- (2,2) -- (2,0)} &
	\hexnhex{(0,0.5) -- (2.5,0.5) -- (2.5,2) -- (1,2) -- (1,1.5)
	    -- (2,1.5) -- (2,1) -- (0.5,1) -- (0.5,2.5) -- (3,2.5)} &
	\hexnhex{(0.5,0) -- ++(0,1.5) -- ++(0.5,0) -- ++(0,1) -- ++(0.5,0) -- ++(0,-2)
	    -- ++(0.5,0) -- ++(0,1) -- ++(0.5,0) -- ++(0,1.5)}}
\end{frame}


\section[Самые сложные задачи]{4. Difficult Problems}

\begin{frame} \frametitle{Карфаген\quad{\it\normalsize (Широкий не значит высокий)}}
\usl{2019-8-1C}{Докажите, что максимальная возможная площадь $n$-угольника, все стороны которого имеют длину 1, меньше, чем максимальная возможная площадь $n+1$-угольника, все стороны которого имеют длину 1.} \vspace{4mm}

Дети же не знают, что максимальную площадь имеют \\
правильные $n$-угольники. \ps

Для каждого многоугольника с $n$ сторонами длины 1 построим многоугольник с $n+1$ сторонами, площадь которого больше.
\end{frame}

\begin{frame} \frametitle{Карфаген\quad{\it\normalsize (Широкий не значит высокий)}}

\begin{center}
\begin{tabular}{ccc}
Если выпуклый & & Если невыпуклый \\
\makecell[c]{
	\definecolor{carthfill}{RGB}{223,223,223}
	\tikz[scale=0.48]{
		\draw[thick] (-1,0) -- ++(60:2) -- (1,0) -- cycle;
		\filldraw[thick,fill=carthfill] (-1,0) ++(195:2) ++(210:2) --
		    ++(30:2) -- ++(15:2) -- ++(0:2) --
		    ++(-15:2) -- ++(-30:2);}
} & \hspace{0.6cm} &
\makecell[c]{
	\includegraphics[scale=1.12]{img/carthage}
}
\end{tabular}\end{center}
\end{frame}

\begin{frame} \frametitle{Сортировка}

Выпишем все числа от одного до десяти — но не в привычном порядке возрастания, а в алфавитном порядке: восемь, два, девять, десять, один, пять, семь, три, четыре, шесть. \medskip

\usl{2020-6-4B}{
Числа от 1 до 10'000'000'000 (десять миллиардов) выписали в алфавитном порядке. Перечислите первые десять из них.
}

(1) 18 (2) 18 миллионов (3) 18 миллионов 18 (4) 18 миллионов 18 тысяч \\
(5) 18 миллионов 18 тысяч 18 (6) \ldots восемь (7) \ldots восемьдесят \\
(8) \ldots 88 (9) \ldots 82 (10) \ldots 89.

\end{frame}


\section[Самые простые задачи]{5. Problems For Peasants}

\begin{frame}
\frametitle{Незакрученный удар}

\usl{2019-4-2A}{Шарик катается
по прямоугольнику, не замедляясь. Когда он подъезжает к краю
прямоугольника, он отскакивает от него и продолжает движение.
В каком положении окажется шарик, будучи запущенным
из центра прямоугольника на рисунке, после того как он проедет
24 клетки по диагонали?}

\def\ballandthick{
	\filldraw (0,0) circle[radius=1.2mm];
	\draw[very thick,->] (0,0) -- (-0.72,0.72);
	\draw[very thick] (-1,-1.5) rectangle (1,1.5);
}

\begin{center} \tikz[scale=1.04]{
	\draw (-1.2,0) node[left]{\small \begin{minipage}{4.2cm}\ \end{minipage}};
	\foreach \x in {-2,...,2} {\draw[color=gray] (-0.5 * \x cm, -1.5) -- (-0.5 * \x cm, 1.5);}
	\foreach \x in {-3,...,3} {\draw[color=gray] (-1, -0.5 * \x cm) -- (1, -0.5 * \x cm);}
	\draw[color=gray, thick, dotted] (0,0) -- (-1,1) -- (-0.5,1.5) -- (0.3,0.7);
	\ballandthick;
\pause
	\draw[color=gray, thick, dotted] (0.3,0.7) -- (1,0) -- (-0.5,-1.5) -- (-1,-1) -- (1,1)
	    -- (0.5,1.5) -- (-1,0) -- (0.5,-1.5) -- (1,-1) -- (0,0);
	\ballandthick; \draw[thick] (-1,0) -- (1,0);
	\draw (1.35,0) node[right]{\small \begin{minipage}{5.6cm} Раз в 6 клеток пересекает \\
	    горизонтальную среднюю линию \end{minipage}};
} \end{center} \end{frame}

\begin{frame} \frametitle{Кирпичей требуют наши сердца}

\usl{2020-4-4B}{
	У Вани есть доски для паркета размером $20 \times 10$ сантиметров, их можно распиливать пополам. Как Ване покрыть этими досками пол квадратной комнаты $1\text{ метр}\times 1\text{ метр}$ так, чтобы не было швов длиной более 30 сантиметров ни в одном из направлений?
}

\begin{center} \vspace{-0.4cm} \tikz[scale=0.4]
{
   \foreach \x / \y in {
	1/2, 1/6, 1/10, 2/1, 2/5, 2/9, 3/4, 3/8, 4/3, 4/7, 5/2, 5/6, 5/10, 6/1, 6/5, 6/9, 7/4, 7/8,
	8/3, 8/7, 9/1, 9/2, 9/6, 9/10
   } 
   \draw[fill opacity=0.6] (\x, \y) rectangle (\x cm + 2 cm, \y cm + 1 cm);
   \foreach \x / \y in {
	1/4, 1/8, 2/3, 2/7, 3/2, 3/6, 4/1, 4/5, 4/9, 5/4, 5/8, 6/3, 6/7, 7/2, 7/6, 8/1, 8/5, 8/9, 9/4, 9/8, 10/3, 10/7
   } 
   \draw[fill opacity=0.6] (\x, \y) rectangle (\x cm + 1 cm, \y cm + 2 cm);
   \foreach \x / \y in {
	1/1, 1/3, 1/7, 3/10, 5/1, 7/10, 10/5, 10/9
   } 
   \draw[fill opacity=0.6] (\x, \y) rectangle (\x cm + 1 cm, \y cm + 1 cm);
   \draw[thick,fill opacity=0.6] (1, 1) rectangle (11 cm, 11 cm);
} 
\end{center}

\end{frame}

\begin{frame} \frametitle{Интереснее, чем кажутся}

\usl{2020-6-6A}{
	Прислонив к зеркалу правый край экрана калькулятора, Серёжа заметил, что отражение некоторых чисел в зеркале тоже можно прочитать как (возможно, другое) число. Например, из 281 получается 185. Можно ли прочитать отражение чисел:
     \[ 180;\quad 205;\quad 12851;\quad 369;\quad 31813\text{?} \]
} \vspace{-0.4cm}

\[ 1 \to 1,\ 2 \to 5,\ 3 \to ?,\ 4 \to ?,\ 5 \to 2,\ 6 \to ?,\ 7 \to ?,\ 8 \to 8,\ 9 \to ?,\ 0 \to 0. \]

\end{frame}


\def\fitem#1#2{\textcolor{hard}{\small (#1)}~~#2 \medskip \\}

\begin{frame} \frametitle{Заключение}

\textcolor{hard}{\bf Направления, которые мы обсудили сегодня}

\fitem{1}{Конструктивные задачи}
\fitem{2}{Задачи, вдохновлённые реальными явлениями}
\fitem{3}{Задачи на приведение {\itshape множества} примеров}
\fitem{4}{Самые сложные задачи для детей с опытом}
\fitem{5}{Самые простые задачи, которые {\itshape решат все} \vspace{10mm}}

\hrule
\begin{center}
	{\LARGE Спасибо за внимание!} \smallskip \\
	{\footnotesize \textcolor{hard}{($*$)\quad}
		Заявки на разбор задач с \url{mathnonstop.ru}
		\phantom{($*$)\quad}}
\end{center} \vspace{2.4mm}
\end{frame}

\end{document}
